\documentclass[12pt,a4paper]{article}
\usepackage{cmap} % Makes the PDF copiable. See http://tex.stackexchange.com/a/64198/25761
\usepackage[T1]{fontenc}
\usepackage[brazil]{babel}
\usepackage[utf8]{inputenc}
\usepackage{amsmath}
\usepackage{amsfonts}
\usepackage{amssymb}
\usepackage{amsthm}
\usepackage{textcomp} % \degree
\usepackage{gensymb} % \degree
\usepackage[usenames,svgnames,dvipsnames]{xcolor}
\usepackage{hyperref}
\usepackage{graphicx}
\usepackage{cancel}
\usepackage[margin=2cm]{geometry}

\hypersetup{
    colorlinks = true,
    allcolors = {blue}
}

% TODO: Consider using exsheets
% http://linorg.usp.br/CTAN/macros/latex/contrib/exsheets/exsheets_en.pdf
%
% http://ctan.org/tex-archive/macros/latex/contrib/exercise/
% Options: answerdelayed,lastexercise,noanswer
\usepackage[answerdelayed,lastexercise]{exercise}

\addto\captionsbrazil{%
\def\listexercisename{Lista de exerc\'icios}%
\def\ExerciseName{Exerc\'icio}%
\def\AnswerName{Solu\c{c}\~ao do exerc\'icio}%
\def\ExerciseListName{Ex.}%
\def\AnswerListName{Solu\c{c}\~ao}%
\def\ExePartName{Parte}%
\def\ArticleOf{de\ }%
}

\renewcommand{\ExerciseHeaderTitle}{(\ExerciseTitle)\ }
\renewcommand{\ExerciseListHeader}{%\ExerciseHeaderDifficulty%
\textbf{%\ExerciseListName\
\ExerciseHeaderNB.\ %
%\ --- \
\ExerciseHeaderTitle}%
%\ExerciseHeaderOrigin
\ignorespaces}
\renewcommand{\AnswerListHeader}{\textbf{\ExerciseHeaderNB.\ (\AnswerListName)\ }}

\newcommand{\fixme}{{\color{red}(...)}}
\newcommand*\diff{\mathop{}\!\mathrm{d}}

\renewcommand{\theenumi}{\alph{enumi}}
\renewcommand\labelenumi{(\theenumi) }

\newcommand*\tipo{Prova II}
\newcommand*\turma{TADS121-01C}
\newcommand*\disciplina{CDI0001}
\newcommand*\eu{Helder G. G. de Lima}
\newcommand*\data{07/10/2016}

\author{\eu}
\title{\tipo - \disciplina}
\date{\data}

\begin{document}
\thispagestyle{empty}
\newgeometry{margin=2cm,bottom=0.5cm}
\begin{center}
\includegraphics[width=9.0cm]{marca} \\
\textbf{\tipo\ (\disciplina / \turma)} \\
Prof. \eu\footnote{
Este é um material de acesso livre distribuído sob os termos da licença \href{https://creativecommons.org/licenses/by-sa/4.0/deed.pt_BR}{Creative Commons BY-SA 4.0}.}
\end{center}

\noindent Nome do(a) aluno(a): \underline{\hspace{9,7cm}} Data: \underline{\data}

%\section*{Instruções}
\begin{center}\fbox{
\begin{minipage}{14cm}

{\footnotesize
\begin{itemize}
\renewcommand{\theenumi}{\Roman{enumi}}
\item Identifique-se em todas as folhas.
\item Mantenha o celular e os demais equipamentos eletrônicos desligados durante a prova.
\item Escolha as questões que irá resolver de modo a totalizar até 10 pontos.
\end{itemize}
}

\end{minipage}
}
\end{center}


%\section*{Questões}
\begin{ExerciseList}
\Exercise%[title={2,0}]
Considere a função $f$ definida por $y = f(x) = \sqrt{4x + 1}$.
\begin{enumerate}
\item \textbf{(1,0)} Calcule $f^\prime(2)$, por meio do limite que define esta derivada.
\item \textbf{(0,5)} Utilize as regras de derivação para obter $f^\prime(a) = \dfrac{ \diff{y} } { {dx} } (a)$, e verifique que para $a = 2$, o valor obtido coincide com o resultado do item anterior.
\item \textbf{(0,5)} Qual é o domínio da função $f^\prime$, isto é, quais são todos os pontos em que $f$ é derivável?
\item \textbf{(2,0)} Utilize a técnica de derivação implícita para calcular $\dfrac{ \diff{x} } { {dy} }$ no ponto $y$ (isto é, obtenha a derivada da função $x=f^{-1}(y)$ definida implicitamente por $y = \sqrt{4x + 1}$).
\end{enumerate}

\Answer
\begin{enumerate}
\item Por definição, tem-se
\begin{align*}
f^\prime(2) & = \lim_{x \to 2} \frac{f(x)-f(2)}{x-2}
   = \lim_{x \to 2} \frac{\sqrt{4x + 1}-\sqrt{4\cdot 2 + 1}}{x-2}
   = \lim_{x \to 2} \frac{\sqrt{4x + 1}-3}{x-2} \\
 & = \lim_{x \to 2} \frac{\sqrt{4x + 1}-3}{x-2} \cdot
                    \frac{\sqrt{4x + 1}+3}{\sqrt{4x + 1}+3}
   = \lim_{x \to 2} \frac{4x + 1 - 9}{(x-2)(\sqrt{4x + 1}+3)} \\
 & = \lim_{x \to 2} \frac{4\cancel{(x - 2)}}{\cancel{(x-2)}(\sqrt{4x + 1}+3)}
   = \lim_{x \to 2} \frac{4}{\sqrt{4x + 1}+3}
   = \frac{4}{\sqrt{4 \cdot 2 + 1}+3}
   = \frac{4}{6}
   = \frac{2}{3}.
\end{align*}

\item Pela regra da cadeia, tem-se
$f^\prime(a) = \dfrac{1}{2\sqrt{4a+1}} \cdot 4 = \dfrac{2}{\sqrt{4a+1}}$, desde que $4a+1 \geq 0$ e $\sqrt{4a+1} \neq 0$, ou seja, para todo $a > -\frac{1}{4}$. Em particular,
\[
f^\prime(2)
= \frac{2}{\sqrt{4\cdot 2+1}}
= \frac{2}{\sqrt{9}}
= \frac{2}{3}.
\]

\item Pelo item anterior, $f^\prime(a)$ está definida para todo $a \in (-\frac{1}{4}, +\infty)$. Então $\operatorname{dom}{f^\prime} = (-\frac{1}{4}, +\infty)$.
\item Se $y = \sqrt{4x + 1}$ e $x=f^{-1}(y)$ então, segue da regra da cadeia que
\[
\dfrac{ \diff } { \diff{y} }(y)
= \dfrac{ \diff } { {dy} } \left(\sqrt{4x + 1}\right)
\Rightarrow
1
= \frac{1}{2\sqrt{4x + 1}}\cdot 4 \cdot  \dfrac{ \diff{x} } { \diff{y} }
\Rightarrow
\frac{2\sqrt{4x + 1}}{4}
= \dfrac{ \diff{x} } { \diff{y} }
\Rightarrow
\frac{2y}{4}
= \dfrac{ \diff{x} } { \diff{y} }.
\]
Portanto, $\dfrac{ \diff{x} } { \diff{y} } = \frac{y}{2}$.
\end{enumerate}

\Exercise[title={2,0}] Seja $g$ a função definida por $g(x) = \ln(12x-3x^2)$. Obtenha a equação da reta tangente ao gráfico de $g$ no ponto $P=(3, g(3))$.
\Answer Sendo $g(x) = \ln(12x-3x^2)$, tem-se
\[
g^\prime(x) = \frac{1}{12x-3x^2} \cdot(12-6x) = \frac{12-6x}{12x-3x^2},
\]
e em particular $g^\prime(3) = \frac{12-6 \cdot 3}{12\cdot 3-3 \cdot 3^2} = \frac{4-6}{12-9} = -\frac{2}{3}$. Além disso, como
\[
g(4) = \ln(12 \cdot 4 - 3 \cdot 4^2) = \ln(36-27) = \ln(9),
\]
o ponto de tangência é $P = (4, \ln(9))$ e a reta que tangencia o gráfico de $g$ neste ponto é $r: y=\ln(9)-\frac{2}{3}(x-3)$, ou seja,
$r: y=-\frac{2}{3}x + 2 - \ln(9)$.

\Exercise%[title={2,0}]
Seja $h(x)=(e^x + 1)^2$.
\begin{enumerate}
\item \textbf{(1,0)} Calcule $h^\prime(x)$, $h^{\prime\prime}(x)$ e $h^{(3)}(x)$.
\item \textbf{(0,5)} Escreva a fórmula geral para a $n$-ésima derivada de $h(x)$.
\item \textbf{(0,5)} Utilize a fórmula do item anterior para calcular $h^{(7)}(0)$.
\end{enumerate}
\Answer
\begin{enumerate}
\item Pela regra da cadeia, tem-se:
\[
h^\prime(x)
= 2(e^x + 1) \cdot e^x
= 2e^{2x} + 2e^x,
\]
\[
h^{\prime\prime}(x)
= 2e^{2x}\cdot 2 + 2e^x
= 4e^{2x} + 2e^x,
\]
\[
h^{(3)}(x)
= 4e^{2x}\cdot 2 + 2e^x
= 8e^{2x} + 2e^x.
\]
\item Considerando que $(e^{2x})^\prime = 2 e^{2x}$, a cada derivação o coeficiente da primeira parcela será duplicado, resultando em
\[
h^{(n)}(x)
= 2^{n}e^{2x} + 2e^x.
\]
\item Para $x = 0$ e $n = 7$, a fórmula anterior fornece:
\[
h^{(n)}(x)
= 2^{7}e^{2 \cdot 0} + 2e^0
= 2^{7} + 2 = 130.
\]
\end{enumerate}

\Exercise[title={2,0}] Utilize uma aproximação por diferenciais (aproximação linear) para estimar o volume de material necessário para fazer uma bola com espessura de $0,1 cm$ e raio interno de $10 cm$.
\Answer Sendo o volume de uma esfera de raio $r$ dado por $V = V(r) = \frac{4}{3} \pi r^3$, tem-se
$\frac{\diff V}{\diff r} = 4 \pi r^2$ e,  consequentemente, $\diff V = 4 \pi r^2 \diff r$. Assim, o volume aproximado é:
\[
\Delta V = V(10,1) - V(10) \approx \diff V = 4 \pi 10^2 \cdot 0,1 = 40 \pi cm^2 = 125,6\ldots cm^2.
\]
\begin{center}
\includegraphics[width=7cm]{img/prova-2-tads-esfera.png}
\end{center}
Observe que $V(10) = \frac{4}{3} \pi 10^3 = \frac{4000}{3} \pi \approx 4188,79$, então o volume exato é bastante próximo do que foi obtido anteriormente:
\[
\Delta V = V(10,1) - V(10) =126,924\ldots cm^2.
\]


\Exercise[title={2,0}] Ao andar por uma calçada, você está se afastando de uma luminária de $2,7m$ de altura a uma velocidade de $1,5 m/s$. Assumindo que você tenha $1,8m$ de altura, utilize o que aprendeu sobre derivadas para determinar com que velocidade o comprimento da sua sombra no chão aumenta conforme você se afasta da luminária.
\Answer Considerando que $x = x(t)$ seja a distância entre você e a luminária no instante de tempo $t$ e que $y = y(t)$ seja o comprimento de sua sombra no mesmo instante, tem-se $\frac{\diff x}{\diff t} = 1,5 m/s = \frac{3}{2} m/s$ e o objetivo é determinar $\frac{\diff y}{\diff t}$.
\begin{center}
\includegraphics[width=8cm]{img/prova-2-tads-sombra.pdf}
\end{center}

Por semelhança de triângulos, conclui-se que
\[
\frac{x+y}{2,7} = \frac{y}{1,8}
\Leftrightarrow
x+y = \frac{2,7 y}{1,8}
\Leftrightarrow
x = \frac{3}{2} y - y = \frac{y}{2}.
\]
Derivando ambos os membros em relação a $t$, resulta que
\[
\frac{\diff x}{\diff t} = \frac{1}{2} \frac{\diff y}{\diff t}
\Rightarrow
\frac{3}{2} = \frac{1}{2} \frac{\diff y}{\diff t}
\Rightarrow
\frac{\diff y}{\diff t} = \frac{6}{2} = 3.
\]
Portanto, o comprimento da sombra aumenta a uma taxa de $3m/s$.
\end{ExerciseList}

\begin{center}
BOA PROVA!
\end{center}

\newpage
\restoregeometry
\section*{Respostas}
\shipoutAnswer
\end{document}
